\section{Disputed Questions after Aquinas}\label{sec:disputed}

On seven questions of the treatise the Catholic schools teach
differently, all within the faith and from the same Scripture and
Fathers. Bonaventure's commentary on the \emph{Sentences}, written in the
years just before the \emph{Prima pars}, was read continuously beside the
\emph{Summa};
the condemnation of two hundred nineteen articles by Stephen Tempier,
bishop of Paris, on 7 March 1277 (the roll itself bears the year 1276, by
the Parisian reckoning that began the year at Easter) fell three years
after Aquinas's death across several of his characteristic positions;
Scotus's \emph{Ordinatio} rebuilt the doctrine of the angels on another
metaphysics; and the commentators of the sixteenth century, Cajetan on
the \emph{Summa} and Suárez in his treatise \emph{De angelis}, carried
the argument forward with full knowledge of both sides.

Aquinas teaches that the angels are altogether immaterial; Bonaventure
and the Franciscan school, that they are composed of a spiritual matter
taken broadly. Aquinas teaches that each angel is its own species;
Bonaventure, that all the angels are of one species as all men are; and
Scotus, that many angels in one species are simply possible. Aquinas and
Cajetan ground the angel's presence in place in the application of his
power; Scotus grounds it in his substance, prior to any operation.
Aquinas holds the angels created together with the visible world as the
more probable opinion, and the Greek Fathers' priority of the angelic
world not erroneous; Suárez reads the \latin{simul} of Lateran~IV as
simultaneity and holds the contrary now rash. Aquinas holds the angels
created in sanctifying grace; the Lombard and Bonaventure, that grace was
added after creation. On the devil's first sin Aquinas places its object
in beatitude sought by natural power, Bonaventure in presumption,
ambition, and envy, Scotus in an inordinate love of self, and Suárez in
the desire of the hypostatic union. On the guardians, Origen weighs
assignment at birth and at baptism, and Jerome, with Aquinas after him,
teaches birth. Two of these questions, the multiplication of angels
within a species and the angel's presence in place, meet the roll of 1277
directly (\S\S\ref{sec:dq-species},~\ref{sec:dq-place}). The witnesses
and moments of the first sin stand side by side in
Appendix~\ref{app:first-sin}.

\subsection{Spiritual Matter in the Angels (\ST{I q.~50 a.~2})}\label{sec:dq-matter}

Aquinas determines that the angels have no matter at all. The operation
follows the substance, understanding is ``an altogether immaterial
operation,'' and so ``it must be that every intellectual substance is
altogether immaterial'' (\ST{I q.~50 a.~2}). The composition the angel
does have, of essence and act of being, is not the composition of matter
and form; the target of the article is Avicebron's \emph{Fount of Life},
which had taught one universal matter for spiritual and corporeal things
alike (\S\ref{sec:q50}).

Bonaventure receives the same authorities and determines the question the
other way. The angel's essence is not simple by the privation of all
composition: it is mutable, and so needs a foundation of possibility;
limited and set in a genus, and so composite in its essence; and so
\latin{illa positio videtur verior esse, scilicet quod in Angelo sit
compositio ex materia et forma} (\emph{In II Sent.} d.~3 p.~1 a.~1 q.~1).
The conclusion is stated with the distinction that governs it:
\latin{Si materia large sumitur extendendo nomen ad omne potentiale
constitutivum, ipsa substantia Angeli composita est ex materia et
forma} -- if matter is taken broadly, as the name for every constitutive
potency, the angelic substance is composed of matter and form. The
authorities who deny matter to spirits, Boethius among them, speak of
matter \latin{appropriate}, in the strict sense of the principle that
undergoes change (ad~1, ad~3).

The two masters are closer than the conclusion suggests and further than
the shared vocabulary suggests. Both hold that the angel is composed of
potency and act. Aquinas concedes the Franciscan name in the wide sense
and refuses it as improper: if every potency is called matter and every
act form, there is matter and form in spiritual substance, \latin{sed
tamen hoc non est proprie dictum secundum communem usum nominum}
(\emph{De spiritualibus creaturis} a.~1). What divides them is whether
the potency in a spiritual creature is matter, with the name and the
doctrine of individuation that the name carries. Bonaventure's spiritual
matter grounds the numerical multiplication of angels within a species;
Aquinas, having refused matter, must and does refuse the multiplication
(\S\ref{sec:dq-species}).

\angeldossier{\ST{I q.~50 a.~2}}
 {Defined: the angelic creature is spiritual and created from nothing
  (Lateran IV, DS~800). Disputed: whether the potency in spiritual
  substance is properly called matter. Thomist position: the angel is
  altogether immaterial (a.~2). Bonaventure and the Franciscan school: a
  spiritual matter in the broad sense (\emph{In II Sent.} d.~3 p.~1 a.~1
  q.~1). Both: the angel is composed of potency and act.}
 {Avicebron, \emph{Fons vitae} (the position examined in a.~2);
  Aristotle, \emph{Physics} I (a.~2); Boethius, \emph{De Trinitate}, as
  Bonaventure reads him, \emph{In II Sent.} d.~3 p.~1 a.~1 q.~1; Aquinas,
  \emph{De spiritualibus creaturis} a.~1.}
 {Avicebron's universal hylomorphism, rejected by Aquinas; the Franciscan
  school's spiritual matter, refused by Aquinas as improper speech
  (\emph{De spir. creat.} a.~1).}

\subsection{Whether Many Angels Share One Species (\ST{I q.~50 a.~4})}\label{sec:dq-species}

Aquinas recites three opinions before his own: that all spiritual
substances, souls included, are of one species; that all the angels are
of one species, but not souls; that all the angels of one hierarchy, or
of one order, are of one species. His determination follows from the
preceding article with the force of an identity: things that agree in
species and differ in number agree in form and are distinguished
materially, and so, the angels having no matter, ``it is impossible for
two angels to be of one species'' (a.~4). Even given matter, he adds,
their diversity of powers would make not merely another species but
another genus.

Bonaventure treats the same question twice. At \emph{In II Sent.} d.~3
p.~1 a.~2 q.~1 (Quaracchi II, 102--104) he sets out two positions. The first, \latin{tot species,
quot individua}, has some probability in bodies but in spirits is
\latin{praesumtio}: Scripture shows many angels ordered to the same
office and the same operation, and no such diversity is made known by
Scripture, the saints, or the angels' functions. The second he calls
\latin{positio sobria et catholica}: \latin{In Angelis, vel in aliquibus
vel in omnibus, est discretio solummodo quoad personalitatem, non quoad
speciem}. At d.~9 q.~1 he commits himself: \latin{videtur magis
theologica et probabilis positio, nisi occurrat manifesta auctoritas in
contrarium, quod omnes Angeli sint eiusdem speciei, sicut et omnes
homines} -- all the angels are of one species, as all men are, with
degrees, offices, and dignities within the species, natural and
gratuitous, as among men.

Scotus states the conclusion against Aquinas without qualification:
\latin{Tenenda est ergo conclusio simpliciter opposita, quod scilicet
simpliciter possibile est plures Angelos esse in eadem specie}
(\emph{Ord.} II d.~3 q.~7). Every quiddity, as quiddity, is communicable
to many: not repugnant to it from perfection, since the divine essence is
communicated to three; not from imperfection, since generable and
corruptible things share one; and only the infinite is incommunicable in
numerical identity. Every created quiddity can be conceived under the
character of a universal without contradiction, which it could not be if
it were of itself \latin{haec}. And God could annihilate this angel and
produce the same species again in another individual, as those who hold
the opposite view must grant in the resurrection of the same man. The
contrary determination rests, in his reading, on taking matter as the
sole principle of numerical distinction.

Between Aquinas and Scotus fell the decree of 1277. Three of its articles
bear on this question: art.~81, \latin{Quod, quia intelligentie non
habent materiam, Deus non posset facere plures ejusdem speciei}; art.~96,
\latin{Quod Deus non potest multiplicare individua sub una specie sine
materia}; and art.~191, \latin{Quod forme non recipiunt divisionem, nisi
per materiam. -- Error, nisi intelligatur de formis eductis de potentia
materie} (\emph{Chartularium Universitatis Parisiensis} I, no.~473). The
roll condemns propositions without naming their authors. Denifle's
apparatus to the document records that these articles \latin{tangere
videntur}, seem to touch, Saint Thomas's teaching on the principle of
individuation; that John of Naples, writing
in the early fourteenth century, reports a current belief that arts.~79,
81, 124, 129, 156, 173, 187, 212, 218, and 219 condemned Thomas's
positions, and argues that none of them derogates from his doctrine; and
that, certainly, \latin{in condemnatione an.~1277 Parisiis non aperte
agitur de S.~Thoma} -- the condemnation does not openly concern Thomas.
Durandus had already joined the articles to this doctrine, recording of
the opinion that God cannot make several angels of one species, \latin{Illa opinio condemnata est
Parisius pluribus vicibus, et de eius condemnatione sunt tres articuli}
-- namely 81, 96, and 191 (Durandus, \emph{In II Sent.} d.~3 q.~3, as
quoted in the Quaracchi scholion to this question, II, 104). Cajetan,
commenting on a.~4, observes that the third of the opinions
Aquinas recites is the one Scotus follows, and presses a doubt against
the proof: the premise that things differing numerically differ
``materially'' is either false or proves nothing, since two souls differ
materially without differing in matter. He does not press the doubt to
the conclusion, and the article's determination stands in his commentary
with the doubt beside it.

\angeldossier{\ST{I q.~50 a.~4}}
 {Disputed: whether several angels can be of one species. Thomist
  position: each angel exhausts its species (a.~4). Bonaventure: personal
  distinction only, in some or in all (\emph{In II Sent.} d.~3 p.~1 a.~2
  q.~1), and all the angels of one species as all men are (d.~9 q.~1).
  Scotus: many angels in one species is simply possible (\emph{Ord.} II
  d.~3 q.~7). Paris 1277 (arts.~81, 96, 191) condemns the proposition
  that immateriality prevents God from multiplying individuals in one
  species.}
 {Dan 7:10 and the angels' common offices (Bonaventure, d.~3 p.~1 a.~2
  q.~1, pp.~102--103); Aristotle, \emph{Metaph.} VII (the opening
  argument at \emph{Ord.} II d.~3 q.~7); \emph{Chartularium Univ.\
  Paris.} I, no.~473, arts.~81, 96, 191; Durandus, \emph{In II Sent.}
  d.~3 q.~3 (in the Quaracchi scholion, II, 104); Cajetan on \ST{I q.~50
  a.~4} (Leonine V).}
 {That two angels cannot be of one species (Aquinas, a.~4), which
  Bonaventure calls presumption and Scotus opposes simpliciter.}

\subsection{The Angel's Presence in Place (\ST{I q.~52 a.~1})}\label{sec:dq-place}

Aquinas distinguishes the senses. A body is in place by the contact of
dimensive quantity; the angel, whose quantity is virtual only, ``is said
to be in a corporeal place by application of the angelic power in any
manner whatever to any place'' (\ST{I q.~52 a.~1}), and is in place as
containing, not as contained. On this follow his positions that an angel
is in only one place at once (a.~2) and that two angels are not in the
same place (a.~3; \S\ref{sec:bodies}).

Scotus takes up the question at \emph{Ord.} II d.~2 q.~5 (Vivès XI;
Cajetan cites the same question as d.~2 q.~6). The first opinion,
identified by his commentator Lychetus as Thomas's at \ST{I q.~52
aa.~1--2}, holds that \latin{Angelus praecise est in loco per
operationem, et nullo modo per essentiam propriam}. Against it Scotus
sets the decree: \latin{istud videtur esse damnatum, sicut quidam
articulus damnatus ab Episcopo Parisiensi et excommunicatus}; and against
the excuse that a Parisian excommunication binds only its own diocese,
\latin{Nec valet dicere, quod excommunicatio non transeat mare, nec
dioecesim; si tamen fuit articulus condemnatus sicut haereticus, ubique
est condemnatus}. Godfrey of Fontaines's variant (\emph{Quodl.} VIII
q.~13), that the angel is in place only by application and not by the
presence of himself, is the same opinion under another word: the
application is either a first act or a second; not a first; not an
immanent second act, since the angel's understanding and willing abstract
from place as his existence does; therefore a transient action, which is
what the first opinion said. Scotus's own arguments turn on the
consequences: an angel at rest in the empyrean, where he works nothing,
would be nowhere; an angel moving a heaven would be in the whole heaven;
he would be in place commensurately; a transient operation is not
formally in the angel; and Damascene joins the angel's being somewhere
and his working there as two, not as one.

The Parisian articles in question are finer than Scotus's use of them.
Art.~204 condemns \latin{Quod substantie separate sunt alicubi per
operationem} -- with the gloss, \latin{Error, si intelligatur, sine
operatione substantiam non esse in loco, nec transire de loco ad locum};
art.~218, \latin{Quod intelligentia, vel angelus, vel anima separata
nusquam est}; art.~219, \latin{Quod substantie separate nusquam sunt
secundum substantiam -- Error, si intelligatur ita, quod substantia non
sit in loco. Si autem intelligatur, quod substantia sit ratio essendi in
loco, verum est}. The glosses condemn the propositions only in the
reading that removes the angel's substance from place apart from
operation; art.~219 is even true in the reading where substance is not
the reason of being in place. John of Naples lists arts.~218 and 219
among the articles believed in his day to strike Thomas, with art.~204
noted beside them as the contrary member.

Cajetan answers for the \emph{Summa} (on \ST{I q.~52 a.~1}, Leonine V).
The relation of a spiritual substance's presence to place always belongs
to an action or a passion, actually or habitually exercised; Scotus's
arguments, which he takes as Capreolus had already gathered them, answer
easily. That the angel should be nowhere in the empyrean is no
absurdity: Aristotle says the blessed beings are not apt to be in place
at all (\emph{De caelo} I), and Gregory Nazianzen holds we speak
improperly when we say the angels are in place, where we should say they
work in place. The gloss of art.~219 turns against the objection itself:
even those who posit a substantial presence in place confess that the
substance is not the \latin{ratio essendi in loco definitive}. And
Aquinas himself, at \emph{Quodl.} I q.~3 a.~1, had already marked the
distinction Cajetan presses: if operation means motion, the proposition
is false; if it means any union of power with place, it is true.

\angeldossier{\ST{I q.~52 aa.~1--3}}
 {Disputed: the ground of the angel's definitive presence. Aquinas and
  Cajetan: application of power, and operation in any sense (a.~1).
  Scotus: a substantial presence naturally prior to operation (\emph{Ord.}
  II d.~2 q.~5). Both: the angel is truly and definitively in place, not
  everywhere (a.~2), and not two in one place (a.~3). Paris 1277
  (arts.~204, 218--219, with their glosses) condemns the angel's being
  nowhere as to substance, in the sense that removes his substance from
  place apart from operation.}
 {Damascene, \emph{De fide orth.} II.3 (a.~1 s.c.); Aristotle,
  \emph{De caelo} I and \emph{Physics} (Cajetan on a.~1);
  \emph{Chartularium Univ.\ Paris.} I, no.~473, arts.~204, 218--219;
  Aquinas, \emph{Quodl.} I q.~3 a.~1 (cited by Cajetan); John of Naples
  on the Thomas articles (in Denifle's apparatus).}
 {That the angel is in place by operation alone and in no way by his
  essence (so Scotus reports and rejects the Thomist position); that the
  angel is nowhere (art.~218) or nowhere as to substance (art.~219),
  censured in the glossed sense.}

\subsection{Before the World, or With It (\ST{I q.~61 a.~3})}\label{sec:dq-time}

Aquinas reports a twofold opinion in the Fathers and grades it himself:
``the more probable one holds that the angels were created at the same
time as corporeal creatures,'' because the angels are part of the one
universe and no part is perfect apart from the whole; ``at the same time
the contrary is not to be deemed erroneous; especially on account of the
opinion of Gregory Nazianzen,'' whose authority in Christian doctrine
Jerome ranks with Athanasius's (a.~3). Jerome, he answers the first
objection, speaks with the Greek Fathers, who all hold the angels'
creation to have preceded the corporeal world's (ad~1).

The Greek line is as Aquinas describes it. Gregory of Nazianzus:
``He first conceived the Heavenly and Angelic Powers\ldots\ Then when His
first creation was in good order, He conceives a second world, material
and visible'' (\emph{Or.} 38.9--10). Basil: ``It appears, indeed, that
even before this world an order of things existed,'' a condition
``outstripping the limits of time,'' in which the Creator perfected
spiritual light and the intellectual and invisible natures, before adding
the sensible world (\emph{Hexaemeron} I.5). Damascene records both
opinions -- some, with Gregory the Theologian, place the angels before
all creation, others after the first heaven; all agree it was before man
-- and declares, ``For myself, I am in harmony with the theologian''
(\emph{De fide orth.} II.3).

Lateran IV's \emph{Firmiter} confesses that God \latin{sua omnipotenti
virtute simul ab initio temporis utramque de nihilo condidit creaturam,
spiritualem et corporalem, angelicam videlicet et mundanam} (DS~800), and
Vatican I's \emph{Dei Filius} repeats the clause word for word, citing
\emph{Firmiter} (DS~3002; \S\ref{sec:faith}). The later dispute concerns
the force of \latin{simul}. Suárez treats it at \emph{De angelis} I.3:
the true and common opinion of the theologians, with the Master and
Aquinas, is that the angels were not created before the corporeal world
but together with it; on the decree he distinguishes. Ferrariensis holds
the simultaneity now pertains to the faith as defined; Vázquez holds it
not defined, and not even definable, since neither Scripture nor
tradition grounds the definition. Suárez reads \latin{simul} as
simultaneity of duration -- the clause stands against the \latin{deinde}
of man's creation that follows it -- yet holds that the Council did not
define the point \latin{ex instituto}, but spoke incidentally, as it
spoke of angelic incorporeity in the same breath; and he notes the
weightiest witness to that reading: Aquinas, writing when the Council had
already been celebrated and its text received among the Decretals, still
held the contrary opinion not erroneous. So, Suárez adds with Bañez, the
contrary cannot now be held without rashness.

\angeldossier{\ST{I q.~61 a.~3}}
 {Defined: God created both the spiritual and the corporeal creature
  from nothing at the beginning of time (Lateran IV, DS~800; Vatican I,
  \emph{Dei Filius}, DS~3002). Common teaching of the Latin Doctors: the
  angels were created together with the corporeal world, more probable
  for Aquinas (a.~3), the true and common opinion for Suárez; the Greek
  Fathers' priority is not erroneous (a.~3). Disputed: what the
  \latin{simul} decides. Ferrariensis: simultaneity defined. Suárez:
  stated by the Council incidentally, its contrary now rash (\emph{De
  angelis} I.3 nn.~12--15). Vázquez: not defined.}
 {Gregory of Nazianzus, \emph{Or.} 38.9--10; Basil, \emph{Hexaemeron}
  I.5; Damascene, \emph{De fide orth.} II.3; Jerome, \emph{In Ep.\ ad
  Tit.} I (a.~3 obj.~1, ad~1); Lateran IV, \emph{Firmiter} (DS~800);
  Suárez, \emph{De angelis} I.3, with Ferrariensis and Vázquez.}
 {The creation of the angels before the corporeal world (the Greek
  Fathers, held not erroneous at a.~3); the definition of simultaneity by
  \emph{Firmiter} (Ferrariensis, against Suárez's reading).}

\subsection{Created in Grace (\ST{I q.~62 a.~3})}\label{sec:dq-grace}

The Master of the Sentences had fixed the older frame. The angels were
created neither blessed nor wretched (II d.~4 c.~1); \latin{Post
creationem namque mox quidam conversi sunt ad Creatorem suum, quidam
aversi}, and the converted were turned by an added grace, \latin{gratia
apposita}, which those who fell never received: \latin{nunquam est
apposita, ut converterentur} (II d.~5 c.~1). On this frame grace follows
creation as a second gift.

Bonaventure keeps the frame and marks the question as one of fact, to
which only arguments of congruity apply, so that either side can be held
probably (\emph{In II Sent.} d.~4 a.~1 q.~2). His own conclusion sides
with the Master and the common run of the doctors: \latin{Probabilius
videtur, Angelos non habuisse gratiam sanctificantem a primo instanti
suae creationis}. The reason that moves him is the angel's unretarded
nature: grace in an angel meets no check, so that an angel made gracious
would be carried wholly into God and could never have sinned; as guilt
fixed the fallen at once, grace would have confirmed them at once.

Aquinas turns the probability the other way: ``Although there are
conflicting opinions on this point\ldots\ it seems more probable, and
more in keeping with the sayings of holy men, that they were created in
sanctifying grace'' (a.~3). The argument is from Augustine's seedlike
forms: as the natural effects were implanted in the creature at its first
fashioning, so grace, which stands to glory as the seedlike form to the
effect, was given the angels from the beginning. His \emph{sed contra} is
Augustine's own sentence: God created the good angels
``with that chaste love by which they cleaved to Him, in one and the same
act creating their nature, and endowing it with grace'' (\emph{De civ.\
Dei} XII.9). To the argument that a gracious angel could not have sinned
-- the argument the older school had pressed -- he answers that every
form inclines after the mode of its subject's nature, and an intellectual
nature inclines freely: the movement of grace does not impose necessity,
and he who has grace can fail to use it (ad~2). The chronology of the
fall turns on this question, for if the devil was
created in grace and could act in the first instant, he must have sinned
at once after it, and if not, an interval may stand between creation and
fall (a.~6; \S\ref{sec:dq-firstsin}).

\angeldossier{\ST{I q.~62 a.~3}}
 {Disputed: whether the angels were created in sanctifying grace from
  the first instant. Aquinas: more probable that they were (a.~3, with
  Augustine, \emph{De civ.\ Dei} XII.9). The Lombard and Bonaventure:
  more probable that grace was added to nature after creation (Lombard,
  II dd.~4--5; Bonaventure, \emph{In II Sent.} d.~4 a.~1 q.~2).}
 {Augustine, \emph{De civ.\ Dei} XII.9 (a.~3 s.c.) and \emph{De Genesi ad
  litteram} VIII.3 (a.~3); 1 John 3:9 (a.~3); Lombard, \emph{Sent.} II
  d.~4 c.~1, d.~5 c.~1; Bonaventure, \emph{In II Sent.} d.~4 a.~1 q.~2.}
 {That grace was not given in the first instant of creation (Lombard's
  frame, common among the doctors before Aquinas); within the same
  question, that a created grace would have necessitated the angel's
  standing (answered at a.~3 ad~2).}

\subsection{The Devil's First Sin (\ST{I q.~63 aa.~2--6})}\label{sec:dq-firstsin}

The faith fixes the frame. Lateran IV confesses that the devil and the
other demons were created good by nature and \latin{ipsi per se facti
sunt mali} -- became evil of themselves (DS~800). Within that frame the
schools agree on the sin's name and divide on its object, its formal
character, and its moment.

That the first sin was pride is the common inheritance. The Lombard:
\latin{Invidiae namque mater est superbia, qua voluerunt se parificare
Deo} (\emph{Sent.} II d.~5 c.~1). Augustine: the bad angels ``have
forsaken Him who supremely is, and have turned to themselves\ldots\ And
this vice, what else is it called than pride? For `pride is the beginning
of sin'{}'' (\emph{De civ.\ Dei} XII.6). Aquinas gives the position its
analysis: a spiritual nature can be affected only by spiritual goods, and
the only sin possible there is the refusal of due subjection to a
superior, which is pride (a.~2). Its object was not equality with God:
the devil knew by natural knowledge that equality was impossible, and no
passion or habit clouded the first act so that he should choose the
impossible; nor did he will to be subject to no one absolutely, which
would be to will his own non-being. He willed likeness: to hold as his
last end a beatitude he could reach by the power of his own nature,
turning from the supernatural beatitude that is had from God's grace --
or, willing the beatitude of grace, to will it from his own power and not
from divine assistance (a.~3). This, he notes, agrees with Anselm, that
the devil ``sought that to which he would have come had he stood fast''
(\emph{De casu diaboli} c.~4). The \emph{De malo} states the same object
in one clause: \latin{eam consequi voluit per virtutem suae naturae; non
tamen sine Deo in naturam operante, sed sine Deo gratiam conferente}
(q.~16 a.~3). Anselm's own chapter runs: \latin{Peccavit ergo volendo
aliquod commodum, quod nec habebat nec tunc velle debuit}; willing what
God willed him not to will, \latin{voluit inordinate similis esse deo},
by a will \latin{nulli subdita} -- subject to no one, which befits God
alone (c.~4).

Bonaventure's account is the most articulated of the older school:
\latin{Primum Angeli peccatum fuit superbia; quod initiatum est in
praesumtione, consummatum in ambitione, confirmatum in invidiae et odii
aversione} (\emph{In II Sent.} d.~5 a.~1 q.~1). The devil presumed as
soon as he beheld his own beauty; presuming, he desired what was wholly
above him -- to preside over the rest by his own authority, \latin{ita
quod nulli subesset; hoc est solius Dei, et hoc est aequiparantiae};
and failing of it, he turned in envy and hatred, in which he was
confirmed.

Scotus dissents on the formal character. Equality with God can be willed
by a simple and inefficacious volition of complacency, \latin{volitio
complacentiae}, which can even have the impossible for its object and
suffices for merit and demerit; Aquinas's impossibility argument touches
only the efficacious volition that seeks means (\emph{Ord.} II d.~6
q.~1). The first inordinate act itself was an inordinate friendship-love
toward himself -- every \latin{velle concupiscentiae} presupposes a
\latin{velle amicitiae} for the one to whom the good is willed -- and the
desire of beatitude for himself followed immoderately (q.~2). Therefore
\latin{primum peccatum ejus non fuit superbia proprie dicta, sed propter
delectationem quam importabat, magis videtur reduci ad luxuriam}, a
spiritual luxuria, as one who delights inordinately in a truth of
geometry; it can still be called presumption, and so pride in the loose
sense, since the appetite for excellence over others follows from it. The
Quaracchi editors, glossing Bonaventure's question, mark exactly this
divergence as Scotus's.

Suárez carries the question into the Christological object. Whether
Lucifer could sin by desiring the hypostatic union for his own nature is
affirmed, given the decree of the Incarnation and its revelation; and
Suárez holds it \latin{valde probabilis} that this was in fact the sin:
\latin{Luciferum de facto peccasse per superbiam, appetendo hypostaticam
unionem, et a principio adversarium Christi fuisse} (\emph{De angelis}
VII.13 n.~13). God \latin{revelavit a principio Christum, ut caput eorum,
ut auctorem gratiae} (VII.30), so that the refusal could be one of faith,
love, and adoration; a positive precept laid on the angels is possible
and, de facto, very probable (VII.28). The tradition that had the devil
fall by envy of the Incarnation is thereby given scholastic form.

The moment of the sin carries its own dispute. That the devil was wicked
in the first instant of his creation was held by some, on the ground that
he refused righteousness as soon as he was made; Augustine reports the
opinion without asserting it, and absolves it of Manicheism (\emph{De
civ.\ Dei} XI.13). Aquinas records its fate: ``this opinion was
reasonably rejected by the masters as erroneous,'' for Isaiah has Lucifer
rise in the morning and Ezekiel place him in the paradise of God (a.~5;
Isa 14:12; Ezek 28:13); the \emph{De malo} adds that the position
\latin{reprobata fuit ab omnibus magistris tunc Parisiis legentibus}
(q.~16 a.~4). His own reason is drawn from the angel's knowledge: in
the first instant the angel turns to the natural knowledge of himself,
where sin cannot be; only afterward can he turn toward or away from what
is above nature (\emph{De malo} q.~16 a.~4). And so he sinned ``at once
after the first instant,'' if he was created in grace and could elicit a
free act in it -- the two premises of the grace question doing the work
(a.~6; \S\ref{sec:dq-grace}). Augustine's own reading of ``the devil
sinneth from the beginning'' (1 John 3:8) is the mean the masters
adopted: not from the beginning of his created existence, but ``from the
beginning of his sin, when by his pride he had once commenced to sin''
(\emph{De civ.\ Dei} XI.15). The witnesses on the object and the moment
are set side by side in Appendix~\ref{app:first-sin}.

\angeldossier{\ST{I q.~63 aa.~2--6}}
 {Defined: the demons were created good and became evil of themselves
  (Lateran IV, DS~800). Common teaching: the first sin was pride in the
  broad sense (Augustine, \emph{De civ.\ Dei} XII.6; Lombard, II d.~5
  c.~1; Bonaventure; Aquinas). Disputed: its object and formal
  character---likeness to God in beatitude sought by natural power
  (Aquinas, a.~3; \emph{De malo} q.~16 a.~3; Anselm, \emph{De casu}
  c.~4); presumption, ambition, and envy (Bonaventure, \emph{In II Sent.}
  d.~5 a.~1 q.~1); inordinate self-love reducible to spiritual luxuria
  rather than pride properly so called (Scotus, \emph{Ord.} II d.~6
  q.~2); desire of the hypostatic union (Suárez, \emph{De angelis}
  VII.13--14, as very probable). Disputed: its moment, the masters having
  rejected a sin in the first instant (a.~5; \emph{De malo} q.~16 a.~4),
  with the interval turning on the grace question (a.~6).}
 {Isa 14:12; Ezek 28:13 (a.~5); 1 John 3:8 (Augustine, \emph{De civ.\
  Dei} XI.15); Ecclus 10:15 (via \emph{De civ.\ Dei} XII.6); Lombard,
  \emph{Sent.} II d.~5; Anselm, \emph{De casu diaboli} c.~4; Dionysius,
  \emph{DN} 4 (a.~4 s.c.; \emph{De malo} q.~16 a.~4 cites \emph{DN} 7);
  Damascene, \emph{De fide orth.} II.4 (a.~7).}
 {That the devil desired equality with God outright (denied by Aquinas
  a.~3, affirmed possible as inefficacious complacency by Scotus, II d.~6
  q.~1); that he sinned in the first instant of creation (rejected by the
  masters as erroneous, a.~5).}

\subsection{Guardians from Birth or from Baptism (\ST{I q.~113 a.~5})}\label{sec:dq-guardians}

The question is as old as the commentary tradition on Matthew 18:10,
``their angels in heaven always see the face of my Father.'' Origen
states it as an open inquiry: at what time do the angels of the little
ones assume the guardianship -- ``from\ldots\ the laver of
regeneration,'' or ``from birth\ldots\ according to the foreknowledge and
predestination of God'' (\emph{Comm.\ in Matt.} XIII.27). For the birth
view he strings the texts: ``He who separated me from my mother's womb''
(Gal 1:15), ``From the womb of my mother thou hast been my protector''
(Ps 70[71]:6), and the like. A second exposition has Christ consign the
believer to a holy angel after regeneration, none attending the
unbelieving but the angels of Satan; a third, that the angel assigned at
birth might himself change as his charge changes, he offers and expressly
leaves to the reader (XIII.28).

Jerome determines it in a sentence that became the school's \emph{sed
contra}: \latin{Magna dignitas animarum, ut unaquaeque habeat ab ortu
nativitatis in custodiam sui angelum delegatum} -- it is the great
dignity of souls that each has an angel delegated for its guardianship
from the birth (\emph{In Matt.} III, on 18:10). Aquinas recites both
opinions from Origen's tract and judges Jerome's the reasonable one.
Benefits conferred on man as a Christian, such as the Eucharist, begin
with baptism; those conferred on him as a rational being begin at birth,
when he receives that nature; and the guardianship of angels is of the
latter kind, as the articles on the guardianship have established (a.~5,
with aa.~1 and 4). In the womb the child is not yet wholly separate from
the mother, and it can be said with some probability that the mother's
angel guards it until birth (ad~3). The Catechism keeps the same width:
``From infancy to death human life is surrounded by their watchful care
and intercession'' (CCC 336; \S\ref{sec:guardians}).

\angeldossier{\ST{I q.~113 a.~5}}
 {Church teaching: human life from infancy to death is surrounded by
  the angels' watchful care (CCC 336; Matt 18:10). Disputed in the
  patristic tradition: assignment of the guardian at birth or at baptism
  (Origen, \emph{Comm.\ in Matt.} XIII.27, weighing both). The schools,
  with Jerome: from birth (a.~5).}
 {Matt 18:10; Origen, \emph{Comm.\ in Matt.} XIII.27--28 (a.~5);
  Jerome, \emph{In Matt.} III on 18:10 (a.~5 s.c.); Gal 1:15, Ps
  70[71]:6 (Origen); Heb 1:14 (a.~5 obj.~1); CCC 336.}
 {Assignment at baptism rather than at birth (the first of Origen's two
  opinions, a.~5 obj.~1); a guardian's presence beginning only after
  regeneration (Origen's second exposition, XIII.28).}
